\newpage
\section{Polares da aeronave}

% Itens 5, 6 e 7 do roteiro: varreduras de alpha em mn 0.85 para os quatro
% casos (fwd/aft, livre/compensado CM = 0), cada polar terminando no CLmax
% da secao critica e com CLmin = -0.5. Marcar o ponto de projeto em todas
% as curvas. PREENCHER apos as rodadas.

% Grafico unico CD x CL com as quatro polares (item 5):
% \begin{figure}[H] ... images/polar_cd_cl.png ... \end{figure}

\begin{table}[H]
\centering
\caption{$C_D$ no $C_L$ de projeto para cada configuração e comparação com o \texttt{designTool}.}\label{tab:cd_projeto}
\renewcommand{\arraystretch}{1.25}
\begin{tabular}{lcc}
\toprule
Configuração & $C_D$ em $C_L = 0{,}5053$ & Desvio do \texttt{designTool} \\
\midrule
CG dianteiro sem deflexões & --- & --- \\
CG traseiro sem deflexões & --- & --- \\
CG dianteiro compensado ($C_M = 0$) & --- & --- \\
CG traseiro compensado ($C_M = 0$) & --- & --- \\
\texttt{designTool} (polar do Lab 02) & 0,02407 & referência \\
\bottomrule
\end{tabular}
\end{table}

\subsection*{Curvas de sustentação}

% Grafico unico CL x alpha dos quatro casos (item 6), ponto de projeto marcado.

\subsection*{Deflexão de profundor}

% Grafico unico CL x delta_e dos quatro casos (item 7), ponto de projeto
% marcado. Discutir se as deflexoes estao em niveis adequados.
